\section{PMF and likelihood as slices of one surface}
\label{sec:slices}

The likelihood is not \emph{computed} from the PMF by applying an operation; it
\emph{is} the PMF, with the roles of argument and parameter exchanged.
``Same functional form'' is not a derivation and not a coincidence --- it is the
definition of the likelihood function.

Fix $n$ and consider the single bivariate function
\begin{equation}
  g(v, q) = \binom{n}{v} q^v (1-q)^{n-v},
  \qquad (v, q) \in \{0,1,\dots,n\} \times [0,1].
\end{equation}
The PMF and the likelihood are two orthogonal slices of this one surface:
\begin{itemize}[nosep]
  \item Fix $q$, let $v$ vary over $\{0,\dots,n\}$: this is the PMF
    $\dens_\Upsilon(v; q, n)$, a slice along the discrete axis.
  \item Fix the observed $v$, let $q$ vary over $[0,1]$: this is
    $\like(q \given v, n)$, a slice along the continuous axis.
\end{itemize}
So the ``computation'' is only: substitute the observed $v$ and read what
remains as a function of $q$. The binomial coefficient and both power terms stay
algebraically identical, which is why the two expressions are typographically
the same. The discrete-versus-continuous distinction is entirely about which
axis was left free: $v$ ranges over a finite set, $q$ over a continuum.

\paragraph{Pointwise equality is not identity of objects.}
The two slices normalize on different axes. Across the discrete axis the PMF
sums to $1$ for every $q$, by the binomial theorem:
\begin{equation}
  \sum_{v=0}^{n} g(v, q) = \bigl(q + (1-q)\bigr)^n = 1.
\end{equation}
Across the continuous axis the likelihood integrates to
\begin{equation}
  \int_0^1 g(v, q)\, dq
  = \binom{n}{v}\, \Betafn(v+1,\, n-v+1)
  = \binom{n}{v}\, \frac{v!\,(n-v)!}{(n+1)!}
  = \frac{1}{n+1},
  \label{eq:lik-integral}
\end{equation}
independent of $v$, and in particular $\neq 1$. Thus $\like$ is not a density in
$q$; this is exactly why maximizing it, rather than integrating it, is the
sensible operation, and why turning it into a density requires the Beta
normalization.

\begin{remark}
That the integral \eqref{eq:lik-integral} equals $\tfrac{1}{n+1}$ for every $v$
is the statement that a uniform prior on $q$ induces a uniform prior-predictive
on the count $v \in \{0,\dots,n\}$ --- the $\mathrm{Beta}(1,1)$--Binomial edge
case.
\end{remark}
